\chapter{System architecture overview}
\label{ch:system-architecture}

This book is about deterministic UI runtimes.
So we start the way systems work starts: with architecture.

Before we zoom into code, this chapter gives you the whole machine:
what subsystems exist, what contracts bind them, and how \emph{time} and \emph{state} move through the system.

You should walk away with one mental model:
	extbf{state marks things dirty, the scheduler orders work, and the renderer commits DOM writes.}

\section{The subsystems}

Relax.js is best understood as a set of cooperating subsystems with explicit contracts:

\begin{itemize}
  \item \textbf{VDOM runtime} (implementation in \texttt{src/*}): general-case rendering and components.
  \item \textbf{Scheduler} (\texttt{src/scheduler.ts} and \texttt{runtime/scheduler.ts}): deterministic work ordering and fairness.
  \item \textbf{Signals/reactive graph} (\texttt{runtime/signals.ts}): fine-grained invalidation and effect scheduling.
  \item \textbf{HRBR runtime} (\texttt{runtime/*}): compiled blocks (template + slots), hydration, fallback reconciliation.
  \item \textbf{Compiler/transform} (\texttt{compiler/*}): optional TSX/JSX transform that emits HRBR blocks or routes to fallback.
\end{itemize}

The important design choice is separation of concerns:
	extbf{reactivity} decides \emph{what} changed, \textbf{scheduling} decides \emph{when} to apply changes, and \textbf{rendering} decides \emph{how} to write to the DOM.

\section{Diagram: end-to-end data flow (authoring \textrightarrow DOM)}

\begin{center}
\begin{tikzpicture}[
  node distance=10mm and 14mm,
  box/.style={draw, rounded corners, align=left, inner sep=6pt, fill=RelaxLight},
  arrow/.style={-Latex, thick}
]
\node[box] (state) {State\\(signals / component state)};
\node[box, right=of state] (render) {Render\\(TSX/JSX or \texttt{h()})};
\node[box, right=of render] (model) {Model\\VNode tree \emph{or} HRBR vnode};
\node[box, right=of model] (commit) {Commit\\DOM writes};
\node[box, below=of commit] (dom) {Real DOM\\(browser/jsdom)};

\draw[arrow] (state) -- node[midway, above] {change} (render);
\draw[arrow] (render) -- node[midway, above] {produces} (model);
\draw[arrow] (model) -- node[midway, above] {mount/patch} (commit);
\draw[arrow] (commit) -- (dom);
\end{tikzpicture}
\end{center}

This entire book is an exploration of the \emph{commit} box: how to make DOM writes correct, deterministic, and efficient.

\section{VDOM flow (general case)}

VDOM is the baseline engine.
It is optimized for flexibility and correctness:

\begin{itemize}
  \item \textbf{Model}: build a VNode tree (plain objects).
  \item \textbf{Mount}: turn it into DOM (store back-references).
  \item \textbf{Patch}: reconcile old/new trees (reuse identity when safe).
  \item \textbf{Destroy}: remove nodes and listeners, release references.
\end{itemize}

\section{HRBR flow (compiler-assisted fast path)}

HRBR is a fast path for stable structure.
It replaces repeated discovery work with precomputed write locations.

\begin{itemize}
  \item The compiler extracts a stable \textbf{template} (HTML).
  \item Dynamic values become \textbf{slots} with precomputed DOM paths.
  \item Updates become direct slot writes scheduled deterministically.
\end{itemize}

\section{Scheduler + signals integration (where determinism comes from)}

Relax leans into deterministic runtime engineering:

\begin{itemize}
  \item \textbf{Signals} compute \emph{what} is dirty (fine-grained invalidation).
  \item \textbf{Schedulers} decide \emph{when} to commit DOM updates (lanes, fairness, budgets).
\end{itemize}

The important contract: reactivity does not "touch the DOM" directly.
It schedules work; the runtime commits batched updates with predictable ordering.

\section{Where to go next}

Chapter~\ref{ch:what-and-why} gives the first-principles renderer contract.
Chapter~\ref{ch:hrbr-motivation} justifies HRBR as an architecture choice.
